\documentclass[a4paper,10pt]{report}\input{../0.Préambule/Préambule_cours_prof}\begin{document}

\definecolor{titlecolor}{rgb}{0.0, 0.48, 0.24}
\newpage
\setcounter{chapter}{12}
\chapter{Les triangles 2}

\section{Somme des angles d'un triangle}

\noindent \begin{minipage}{11.2cm}
Découper un triangle quelconque et réaliser le pliage ci contre de façon à ramener les sommets du triangle pour former un rectangle.

\medskip
\noindent Que constatez vous ?

\noindent\grandscarreaux{\linewidth}{1.6cm}
\end{minipage}
\hspace{1mm}
\begin{minipage}{6cm}
\begin{center}
\includegraphics[scale=0.5]{../40.Image/13.Triangles_2_1.png}
\end{center}
\end{minipage}

\vspace{-0.5mm}
\noindent\grandscarreaux{\linewidth}{2.4cm}

\begin{prop}
La somme des mesures des angles d'un triangle est égale à $180$\degre .
\end{prop}

\begin{ex}
$ABC$ est un triangle tel que $\widehat{ABC}= 80$\degre et $\widehat{BAC}= 40$\degre. Calculer $\widehat{BCA}$.

\medskip\medskip
\noindent \textbf{On sait que :} \hspace{6.7mm}$\bullet$ \dotfill \medskip \medskip%$ABC$ est un triangle.

\hspace{2cm} $\bullet$ \dotfill \medskip \medskip%$\widehat{ABC}= 80$\degre.

\hspace{2cm} $\bullet$ \dotfill \medskip %$\widehat{BAC}= 40$\degre.


\medskip
\noindent \textbf{Propriété :} \dotfill \medskip \medskip

\dotfill \medskip%La somme des mesures des angles d'un triangle est égale à $180$\degre .

\medskip
\noindent \textbf{Conclusion :} 

\medskip
\noindent\grandscarreaux{\linewidth}{3.2cm}

%\vspace{-10mm} 
%\begin{eqnarray*}
%\widehat{ABC} + \widehat{BAC} + \widehat{BCA} &=& 180\\
%80 + 40 + \widehat{BCA} &=& 180 \\
%\widehat{BCA} &=& 180 - 80 -60 \quad = \quad 60
%\end{eqnarray*} 
%La mesure de l'angle $\widehat{BCA}$ est de $60$\degre .
\end{ex}

\section{Triangles particuliers}

\subsection{Triangle isocèle}

\begin{prop}
Dans un triangle isocèle, les deux angles à la base ont la même mesure.
\end{prop}

\begin{ex}
$ABC$ est un triangle isocèle en $A$ tel que $\widehat{C} = 80$\degre . Calculons $\widehat{A}$ et $\widehat{B}$.

\medskip
\noindent \begin{minipage}{12cm}

\medskip
\noindent \textbf{On sait que :} \hspace{6.7mm}$\bullet$ \dotfill \medskip \medskip%$ABC$ est isocèle en $A$

\hspace{2.5cm} $\bullet$ \dotfill \medskip %$\widehat{C}= 80$\degre.

\medskip
\noindent \textbf{Propriété :} \dotfill \medskip \medskip

\dotfill \medskip%Dans un triangle isocèle, les deux angles à la base ont la même mesure. 

\medskip
\noindent \textbf{Conclusion :} 

\medskip
\noindent\grandscarreaux{\linewidth}{1.6cm}%$\widehat{B}$ et $\widehat{C}$ ont la même mesure. Donc : $$\widehat{B} = \widehat{C} = 80$$
%La mesure de l'angle $\widehat{B}$ est de $80$\degre .
\end{minipage}
\hspace{3mm}
\begin{minipage}{4cm}
\begin{tikzpicture}[line cap=round,line join=round,>=triangle 45,x=1.0cm,y=1.0cm,scale= 0.7]
\clip(-1.8,-0.8) rectangle (4.8,4.);
\draw [shift={(3.,3.)},line width=0.8pt,color=red,fill=red,fill opacity=0.1] (0,0) -- (-143.13010235415598:0.6) arc (-143.13010235415598:-71.56505117707799:0.6) -- cycle;
\draw [line width=1.2pt] (-1.,0.)-- (4.,0.);
\draw [line width=1.2pt] (1.46,0.1) -- (1.46,-0.1);
\draw [line width=1.2pt] (1.54,0.1) -- (1.54,-0.1);
\draw [line width=1.2pt] (4.,0.)-- (3.,3.);
\draw [line width=1.2pt] (3.,3.)-- (-1.,0.);
\draw [line width=1.2pt] (1.092,1.444) -- (0.972,1.604);
\draw [line width=1.2pt] (1.028,1.396) -- (0.908,1.556);
\begin{scriptsize}
\draw [color=blue] (-1.,0.)-- ++(-3.5pt,0 pt) -- ++(7.0pt,0 pt) ++(-3.5pt,-3.5pt) -- ++(0 pt,7.0pt);
\draw[color=blue] (-1.46,-0.27) node {\large{$A$}};
\draw [color=blue] (4.,0.)-- ++(-3.5pt,0 pt) -- ++(7.0pt,0 pt) ++(-3.5pt,-3.5pt) -- ++(0 pt,7.0pt);
\draw[color=blue] (4.32,-0.23) node {\large{$B$}};
\draw [color=blue] (3.,3.)-- ++(-3.5pt,0 pt) -- ++(7.0pt,0 pt) ++(-3.5pt,-3.5pt) -- ++(0 pt,7.0pt);
\draw[color=blue] (3.3,3.49) node {\large{$C$}};
\end{scriptsize}
\end{tikzpicture}
\end{minipage} 

\vspace{-0.5mm}
\noindent\grandscarreaux{\linewidth}{4cm}
%\medskip
%\noindent On sait que $\widehat{A} + \widehat{B} + \widehat{C} = 180$\degre
%\begin{eqnarray*}
%\widehat{A} + \widehat{B} + \widehat{C} &=& 180\\
%\widehat{A} + 80 + 80 &=& 180 \\
%\widehat{A} + 160 &=& 180 \\
%\widehat{A} &=& 180 - 160 \\
%&=&20
%\end{eqnarray*}
%La mesure de l'angle $\widehat{A}$ est de $20$\degre .
\end{ex}

\begin{ex}
Le triangle $HPK$ est un triangle isocèle en $P$ tel que : $\widehat{P}= 70$\degre. Calculons $\widehat{H}$ et $\widehat{K}$.

%\vspace{-11mm}
\medskip

\noindent \begin{minipage}{12cm}

\noindent\grandscarreaux{\linewidth}{3.2cm}
%\vspace{5mm}
%On sait que $\widehat{H} + \widehat{K} + \widehat{P} = 180$\degre
%\begin{eqnarray*}
%\widehat{H} + \widehat{K} + \widehat{P} &=& 180\\
%\widehat{H} + \widehat{K} + 70 &=& 180 \\
%\widehat{H} + \widehat{K} &=& 180 - 70 \\
%&=&110
%\end{eqnarray*} 
\end{minipage}
\begin{minipage}{4cm}
\begin{center}
\begin{tikzpicture}[line cap=round,line join=round,>=triangle 45,x=1.0cm,y=1.0cm,scale=0.6]
\clip(-0.8,-0.8) rectangle (4.8,6);
\draw [shift={(2.,5.)},line width=0.8pt,color=red,fill=red,fill opacity=0.1] (0,0) -- (-111.80140948635182:0.6) arc (-111.80140948635182:-68.19859051364821:0.6) -- cycle;
\draw [line width=1.2pt] (2.,5.)-- (0.,0.);
\draw [line width=1.2pt] (1.1077032961426903,2.5) -- (0.9220079579656385,2.574278135270821);
\draw [line width=1.2pt] (1.0779920420343616,2.425721864729179) -- (0.8922967038573097,2.5);
\draw [line width=1.2pt] (0.,0.)-- (4.,0.);
\draw [line width=1.2pt] (4.,0.)-- (2.,5.);
\draw [line width=1.2pt] (2.9220079579656386,2.425721864729179) -- (3.1077032961426903,2.5);
\draw [line width=1.2pt] (2.8922967038573097,2.5) -- (3.0779920420343614,2.574278135270821);
\begin{scriptsize}
\draw [color=blue] (2.,5.)-- ++(-3.5pt,0 pt) -- ++(7.0pt,0 pt) ++(-3.5pt,-3.5pt) -- ++(0 pt,7.0pt);
\draw[color=blue] (2.,5.67) node {\large{$P$}};
\draw [color=blue] (0.,0.)-- ++(-3.5pt,0 pt) -- ++(7.0pt,0 pt) ++(-3.5pt,-3.5pt) -- ++(0 pt,7.0pt);
\draw[color=blue] (-0.52,-0.31) node {\large{$K$}};
\draw [color=blue] (4.,0.)-- ++(-3.5pt,0 pt) -- ++(7.0pt,0 pt) ++(-3.5pt,-3.5pt) -- ++(0 pt,7.0pt);
\draw[color=blue] (4.4,-0.15) node {\large{$H$}};
\end{scriptsize}
\end{tikzpicture}
\end{center}
\end{minipage} 

\medskip
\noindent \textbf{On sait que :} \hspace{6.7mm}$\bullet$ \dotfill \medskip %$HPK$ est isocèle en $P$

\medskip
\noindent \textbf{Propriété :} \dotfill \medskip \medskip%Dans un triangle isocèle, les deux angles à la base ont la même mesure. 

\dotfill 

\medskip
\noindent \textbf{Conclusion :} 

\medskip
\noindent\grandscarreaux{\linewidth}{4cm}%$\widehat{H}$ et $\widehat{K}$ ont la même mesure. Donc : $$\widehat{H} = \widehat{K} = \dfrac{110}{2} = 55$$
%La mesure des angles $\widehat{H}$ et $\widehat{K}$ est de $55$\degre .
\end{ex}

\subsection{Triangle équilatéral}

\vspace{-10mm}
\noindent \begin{minipage}{11cm}

\vspace{5mm}
\begin{prop}
Dans un triangle équilatéral, les trois angles sont tous égaux à $60$\degre
\end{prop}
\end{minipage}
\hspace{10mm}
\begin{minipage}{4cm}
\begin{tikzpicture}[line cap=round,line join=round,>=triangle 45,x=1.0cm,y=1.0cm,scale=0.7]
\clip(-1.8,-0.8) rectangle (4.8,5.2);
\draw [shift={(1.48,4.3)},line width=0.8pt,color=red,fill=red,fill opacity=0.1] (0,0) -- (-119.97394866023606:0.6) arc (-119.97394866023606:-59.62771364235506:0.6) -- cycle;
\draw [shift={(-1.,0.)},line width=0.8pt,color=red,fill=red,fill opacity=0.1] (0,0) -- (0.:0.6) arc (0.:60.02605133976395:0.6) -- cycle;
\draw [shift={(4.,0.)},line width=0.8pt,color=red,fill=red,fill opacity=0.1] (0,0) -- (120.37228635764494:0.6) arc (120.37228635764494:180.:0.6) -- cycle;
\draw [line width=1.2pt] (-1.,0.)-- (4.,0.);
\draw [line width=1.2pt] (1.46,0.1) -- (1.46,-0.1);
\draw [line width=1.2pt] (1.54,0.1) -- (1.54,-0.1);
\draw [line width=1.2pt] (4.,0.)-- (1.48,4.3);
\draw [line width=1.2pt] (2.673948827084454,2.064928015442917) -- (2.846500493874763,2.1660513178409593);
\draw [line width=1.2pt] (2.633499506125237,2.1339486821590405) -- (2.8060511729155464,2.235071984557083);
\draw [line width=1.2pt] (1.48,4.3)-- (-1.,0.);
\draw [line width=1.2pt] (0.34660951280711516,2.1346894879566505) -- (0.17335898179026119,2.234610724450091);
\draw [line width=1.2pt] (0.3066410182097388,2.065389275549909) -- (0.13339048719288485,2.16531051204335);
\begin{scriptsize}
\draw [color=blue] (-1.,0.)-- ++(-3.5pt,0 pt) -- ++(7.0pt,0 pt) ++(-3.5pt,-3.5pt) -- ++(0 pt,7.0pt);
\draw[color=blue] (-1.46,-0.27) node {\large{$L$}};
\draw [color=blue] (4.,0.)-- ++(-3.5pt,0 pt) -- ++(7.0pt,0 pt) ++(-3.5pt,-3.5pt) -- ++(0 pt,7.0pt);
\draw[color=blue] (4.32,-0.23) node {\large{$C$}};
\draw [color=blue] (1.48,4.3)-- ++(-3.5pt,0 pt) -- ++(7.0pt,0 pt) ++(-3.5pt,-3.5pt) -- ++(0 pt,7.0pt);
\draw[color=blue] (1.78,4.79) node {\large{$A$}};
\end{scriptsize}
\end{tikzpicture}
\end{minipage}

\vspace{-5mm}
\subsection{Triangle rectangle}

\begin{prop}
Dans un triangle rectangle, la somme des mesures des deux angles aigus est égale à $90$\degre
\end{prop}

\begin{ex}
$IJK$ est un triangle rectangle en $I$ tel que $\widehat{J} = 48$\degre . Calculons $\widehat{K}$.

\medskip
\noindent \begin{minipage}{12cm}

\medskip
\noindent \textbf{On sait que :} \hspace{6.7mm}$\bullet$ \dotfill \medskip \medskip%$IJK$ est rectangle en $I$

\hspace{2.5cm}$\bullet$ \dotfill \medskip 

\medskip
\noindent \textbf{Propriété :} \dotfill \medskip \medskip%Dans un triangle rectangle, la somme des mesures des deux angles aigus est égale à $90$\degre

\dotfill \medskip

\medskip
\noindent \textbf{Conclusion :}

%\vspace{-10mm} 
%\begin{eqnarray*}
%\widehat{J} + \widehat{K} &=& 90 \\
%48 + \widehat{K} &=& 90 \\
%\widehat{K} &=& 90 -48 \quad = \quad 42
%\end{eqnarray*} 
%La mesure de l'angle $\widehat{K}$ est de $42$\degre .
\end{minipage}
\begin{minipage}{4cm}
\begin{tikzpicture}[line cap=round,line join=round,>=triangle 45,x=1.0cm,y=1.0cm,scale=0.7]
\clip(-1.8,-0.8) rectangle (4.8,3.8);
\draw [shift={(-1.,2.8)},line width=0.8pt,color=red,fill=red,fill opacity=0.1] (0,0) -- (-90.:0.6) arc (-90.:-29.248826336546973:0.6) -- cycle;
\draw[line width=0.8pt,color=red,fill=red,fill opacity=0.1] (-0.5757359312880715,0.) -- (-0.5757359312880714,0.42426406871192845) -- (-1.,0.42426406871192845) -- (-1.,0.) -- cycle; 
\draw [line width=1.2pt] (-1.,0.)-- (4.,0.);
\draw [line width=1.2pt] (4.,0.)-- (-1.,2.8);
\draw [line width=1.2pt] (-1.,2.8)-- (-1.,0.);
\begin{scriptsize}
\draw [color=blue] (-1.,0.)-- ++(-3.5pt,0 pt) -- ++(7.0pt,0 pt) ++(-3.5pt,-3.5pt) -- ++(0 pt,7.0pt);
\draw[color=blue] (-1.46,-0.27) node {\large{$I$}};
\draw [color=blue] (4.,0.)-- ++(-3.5pt,0 pt) -- ++(7.0pt,0 pt) ++(-3.5pt,-3.5pt) -- ++(0 pt,7.0pt);
\draw[color=blue] (4.32,-0.23) node {\large{$K$}};
\draw [color=blue] (-1.,2.8)-- ++(-3.5pt,0 pt) -- ++(7.0pt,0 pt) ++(-3.5pt,-3.5pt) -- ++(0 pt,7.0pt);
\draw[color=blue] (-0.7,3.29) node {\large{$J$}};
\end{scriptsize}
\end{tikzpicture}
\end{minipage} 

\noindent\grandscarreaux{\linewidth}{4.8cm}
\end{ex}

\end{document}